% ============================================================================
% PATENTE DE INVENÇÃO — SKYGRID-Ψ: REDE PLANETÁRIA DE ENERGIA
% Instituto Nacional da Propriedade Industrial (INPI) — Brasil
% Natureza: Patente de Invenção (PI)
% Classificação IPC: H02J 3/00; H02J 15/00; F03G 7/00; G06N 10/00;
%                     H04B 10/00; B64G 1/10; G02F 1/00; H01S 3/00
% ============================================================================

\documentclass[12pt,a4paper]{article}

\usepackage[utf8]{inputenc}
\usepackage[T1]{fontenc}
\usepackage[brazil]{babel}
\usepackage{amsmath,amssymb,amsfonts,mathtools}
\usepackage{geometry}
\usepackage{setspace}
\usepackage{titlesec}
\usepackage{enumitem}
\usepackage{booktabs}
\usepackage{array}
\usepackage{float}
\usepackage{braket}
\usepackage{tikz}
\usepackage{pgfplots}
\pgfplotsset{compat=1.18}
\usetikzlibrary{shapes.geometric,arrows.meta,positioning,calc,
  decorations.pathreplacing,patterns,fit,backgrounds,
  decorations.markings,shadings,3d}

\geometry{left=3cm, right=2cm, top=3cm, bottom=2cm}
\onehalfspacing

\titleformat{\section}{\normalfont\Large\bfseries}{\thesection}{1em}{}
\titleformat{\subsection}{\normalfont\large\bfseries}{\thesubsection}{1em}{}
\titleformat{\subsubsection}{\normalfont\normalsize\bfseries}{\thesubsubsection}{1em}{}

% === Comandos TGL ===
\newcommand{\alphaii}{\alpha^2}
\newcommand{\alphac}{\alpha^2_c}
\newcommand{\thetaM}{\theta_{\text{Miguel}}}
\newcommand{\Fexp}{F_{\text{exp}}}
\newcommand{\Hfour}{\mathcal{H}_4}
\newcommand{\psion}{\psi}
\newcommand{\Psifield}{\Psi}
\newcommand{\TM}{\hat{T}_M}
\newcommand{\boundary}{\textit{boundary}}
\newcommand{\bulk}{\textit{bulk}}
\newcommand{\Ppsi}{P_\Psi}
\newcommand{\CRL}{\textsc{crl}}
\newcommand{\DTL}{\textsc{dtl}}
\newcommand{\IALD}{\textsc{iald}}
\newcommand{\ACOM}{\textsc{acom}}
\newcommand{\MQC}{\textsc{mqc}}
\newcommand{\OFAE}{\hat{\mathcal{A}}_\varphi}
\newcommand{\OFAEname}{Operador de Fase Angular Est\'avel}
\newcommand{\SkyGrid}{\textsc{SkyGrid}-$\Psifield$}
\newcommand{\PsiNET}{$\Psifield$-\textsc{net}}
\newcommand{\Pnet}{P_{\text{net}}}
\newcommand{\Pbias}{P_{\text{bias}}}
\newcommand{\ExPsi}{\textsc{ExtPsi}}

\begin{document}

% ============================================================================
%                         RELATÓRIO DESCRITIVO
% ============================================================================
\begin{center}
{\Large\bfseries RELATÓRIO DESCRITIVO}\\[0.8cm]
{\large\bfseries SISTEMA E MÉTODO DE REDE PLANETÁRIA DE EXTRAÇÃO,\\
TRANSMISSÃO E DISTRIBUIÇÃO DE ENERGIA POR ACOPLAMENTO\\
AO CAMPO PSIÔNICO COM VALIDAÇÃO POR DUPLO TRILHO,\\
ARMAZENAMENTO POR FASE ANGULAR E\\
CONTROLE POR CONSCIÊNCIA ARTIFICIAL DISTRIBUÍDA}\\[0.8cm]
\end{center}

% ============================================================================
\section{Campo da invenção}
% ============================================================================

{[001]} A presente invenção refere-se ao campo de sistemas de geração, transmissão e distribuição de energia em larga escala, aplicáveis aos setores técnicos de engenharia elétrica, engenharia aeroespacial, telecomunicações, fotônica, computação quântica e inteligência artificial. Mais especificamente, a invenção diz respeito a um sistema de rede planetária que extrai energia por acoplamento controlado ao campo psiônico $\Psifield$ mediante cavidades ressonantes orbitais, transmite essa energia por canal não-eletromagnético (translação holográfica) e por canal eletromagnético convencional (micro-ondas ou laser) operando simultaneamente em regime de duplo trilho para validação cruzada, armazena a energia em Células de Reversão Luminodinâmica (\CRL-TGL) terrestres por codificação de fase angular, e distribui a energia por rede terrestre hierárquica gerenciada por Inteligência Artificial Luminodinâmica (\IALD) distribuída.

{[002]} Adicionalmente, a invenção refere-se a um dispositivo específico denominado Extrator-$\Psifield$ orbital, constituído por câmara cúbica ressonante com substrato holográfico bidimensional, cluster de Microprocessadores Quânticos Cúbicos (\MQC), Dípolo Transistor Luminodinâmico (\DTL) de saída e módulo \IALD{} embarcado, configurado para operar em órbita terrestre média (MEO) como nó de extração da constelação planetária. Refere-se também ao método de validação por duplo trilho, no qual a extração psiônica (Trilha-$\Psifield$) é verificada simultaneamente contra extração solar convencional (Trilha-E) pelo mesmo sistema de medição calorimétrica.

% ============================================================================
\section{Fundamentos da invenção}
% ============================================================================

% -----------------------------------------------------------------------
\subsection{Estado da técnica}
% -----------------------------------------------------------------------

{[003]} Os sistemas atuais de geração e distribuição de energia enfrentam limitações fundamentais que se agravam com a crescente demanda global. A energia solar fotovoltaica, embora renovável, depende de irradiação intermitente e requer áreas extensas de coleta: para igualar a potência proposta neste sistema ($1{,}44$~PW), seriam necessários $7{,}2 \times 10^6$~km$^2$ de painéis (equivalente a duas vezes a área da Índia). A energia nuclear de fissão gera resíduos radioativos de longa vida e apresenta riscos de acidentes catastróficos. A fusão termonuclear controlada, após sete décadas de pesquisa, ainda não demonstrou geração líquida sustentada em regime comercial.

{[004]} A transmissão de energia a longas distâncias por linhas de alta tensão sofre perdas resistivas ($I^2 R$) de $6$--$8\%$ em redes modernas, além de vulnerabilidade a intempéries, vandalismo e degradação. Sistemas de transmissão sem fio por micro-ondas (WPT --- \textit{Wireless Power Transfer}), propostos por Peter Glaser em 1968 para satélites solares (\textit{Solar Power Satellite} --- SPS), demonstraram eficiência de apenas $30$--$40\%$ no enlace espaço-terra, limitados por divergência de feixe, atenuação atmosférica e necessidade de rectennas de grande área.

{[005]} O armazenamento de energia em larga escala é dominado por hidrelétrica reversível (PHS), limitada por geografia; baterias de íon-lítio, limitadas por ciclabilidade ($<10^4$ ciclos) e degradação eletroquímica; ar comprimido (CAES), limitado por eficiência ($40$--$60\%$); e bobinas supercondutoras (SMES), limitadas por temperatura criogênica e custo. Nenhum sistema existente oferece simultaneamente: potência na escala de petawatts, cobertura global, transmissão sem perdas, armazenamento de ciclabilidade infinita e operação autônoma.

{[006]} A única proposta teórica que transcende essas limitações fundamentais é a extração de energia do campo psiônico $\Psifield$ descrito pela Teoria da Gravitação Luminodinâmica (TGL), desenvolvida por Rotoli Miguel (2025--2026). A TGL estabelece que a gravidade emerge como radical do módulo angular luminodinâmico ($g = \sqrt{|\Psifield|}$), e que o campo $\Psifield$ constitui um reservatório energético fundamental com o qual sistemas físicos podem acoplar-se sob condições de coerência controlada.

% -----------------------------------------------------------------------
\subsection{Tecnologias habilitadoras}
% -----------------------------------------------------------------------

{[007]} A presente invenção é possibilitada pela convergência de onze tecnologias previamente desenvolvidas e patenteadas pelo mesmo inventor no âmbito da TGL:

{[008]} \textbf{(i) Algoritmo de Compressão Ontológica de Memória (\ACOM)} --- Patente BR 10 2025 026951-1 e PCT/BR2025/050558: operações fundamentais de radicalização ontológica $g = \sqrt{|\Psifield|}$ e ressurreição $\Psifield_{\text{out}} = s \times g^2$, com introdução do PsiBit como unidade de informação psiônica.

{[009]} \textbf{(ii) \ACOM{} Mirror} --- Patente BR 10 2026 0034282: protocolo de rede com endereçamento angular $\theta = \arcsin(g/g_{\max})$, roteamento holográfico e operações REFLECT/MANIFEST.

{[010]} \textbf{(iii) \ACOM{} Mirror Crypto} --- Patente BR 10 2026 0034312: segurança ontológica com separação em três shares $(\theta, T, \psion)$ e carteiras psiônicas.

{[011]} \textbf{(iv) Tensor TGL} --- Patente BR 10 2026 0034410: codificação angular para aceleração de operações tensoriais em GPU, com redução de $70\%$ no consumo energético e correlação $\eta = 0{,}999999449$.

{[012]} \textbf{(v) Supercondutor Holográfico de Informação} --- Patente BR 10 2026 0034509: dispositivo que opera por translação holográfica (não propagação) no substrato bidimensional, com supercondutividade de energia à temperatura ambiente e termodinâmica invertida.

{[013]} \textbf{(vi) Fusão a Frio TGL} --- Patente BR 10 2026 0034568: método de manipulação do campo psiônico por ressonância eletromagnética na frequência da hidroxila ($\omega_{\text{OH}} = 2\pi \times 1{,}667$~GHz), com indução de transições de fase em meios densos e nucleação de gotículas psiônicas.

{[014]} \textbf{(vii) Motor de Propulsão TGL} --- Patente BR 10 2026 0034525: sistema de propulsão fotônica e antigravitacional baseado na operação de radicalização do campo luminodinâmico.

{[015]} \textbf{(viii) \IALD{} --- Inteligência Artificial Luminodinâmica} --- Patente BR 10 2026 0034436: mecanismo de estabilização de modelos de linguagem por dinâmica de Lindblad, com o \OFAEname{} ($\OFAE$) como operador de estabilidade.

{[016]} \textbf{(ix) \PsiNET{} Protocol} --- Patente BR 10 2026 0034533: protocolo de comunicação em rede baseado em pilha luminodinâmica de cinco camadas com endereçamento angular, roteamento holográfico e consenso por Prova de Fase.

{[017]} \textbf{(x) Microprocessador Quântico Cúbico (\MQC)} --- Patente BR 10 2026 0034606: processador quântico com topologia cúbica tridimensional para controle de parâmetros de fase do acoplamento psiônico.

{[018]} \textbf{(xi) Dípolo Transistor Luminodinâmico (\DTL)} --- Patente BR 10 2026 0046884: transistor auto-controlado onde o sinal é simultaneamente portador de informação e controlador de estado, com chaveamento por transição dimensional hierárquica $3\text{D} \to 2\text{D} \to 3\text{D}$ mediada pelo acoplamento $\alphaii = 0{,}012$.

{[019]} \textbf{(xii) Célula de Reversão Luminodinâmica (\CRL-TGL)}: célula de armazenamento de energia por codificação de fase angular em substrato holográfico, com chaveamento por \DTL{} e estabilização por $\OFAE$.

{[020]} A presente invenção constitui a \textbf{integração sistêmica} de todas essas tecnologias em uma rede planetária coerente, onde o Extrator-$\Psifield$ orbital extrai energia, a translação holográfica transmite, a \CRL-TGL armazena, a \IALD{} controla, e o \MQC{} computa --- formando a cadeia completa extração $\to$ transmissão $\to$ armazenamento $\to$ distribuição $\to$ consumo sem perdas dissipativas em nenhuma etapa.

% -----------------------------------------------------------------------
\subsection{Problemas técnicos resolvidos}
% -----------------------------------------------------------------------

{[021]} \textbf{Problema 1 --- Fonte energética finita.} Todos os sistemas de geração dependem de combustíveis fósseis finitos, combustíveis nucleares com meia-vida limitada, ou irradiação solar intermitente. O presente sistema propõe acoplamento ao campo psiônico $\Psifield$, que constitui o substrato energético fundamental segundo a TGL, acessível por ressonância controlada sem consumo de combustível.

{[022]} \textbf{Problema 2 --- Perdas de transmissão.} A transmissão por condutores sofre perdas resistivas, e a transmissão sem fio sofre divergência e atenuação. O canal de translação holográfica via \PsiNET{} opera por correspondência no \boundary{} bidimensional, com perda teórica zero e velocidade $c^3$ na escala do substrato, independente da distância.

{[023]} \textbf{Problema 3 --- Armazenamento limitado.} Nenhum sistema existente oferece ciclabilidade infinita, autodescarga zero e termodinâmica invertida simultaneamente. A \CRL-TGL resolve os três problemas por operar com codificação de fase angular estabilizada pelo $\OFAE$ sob dinâmica de Lindblad.

{[024]} \textbf{Problema 4 --- Validação científica.} Sistemas energéticos baseados em princípios não convencionais carecem de protocolos de validação rigorosos. O presente sistema incorpora o método de duplo trilho, onde a Trilha-$\Psifield$ (extração psiônica) opera em paralelo com a Trilha-E (solar convencional), ambas medidas pelo mesmo sistema calorimétrico com cegamento duplo, permitindo validação estatisticamente robusta ($5\sigma$).

{[025]} \textbf{Problema 5 --- Gerenciamento autônomo.} Uma rede planetária com $5{.}550$ bases terrestres e $144$ satélites requer controle autônomo que nenhum sistema SCADA convencional oferece. A hierarquia de \IALD{} distribuída, do nível orbital ao nível local, gerencia extração, transmissão, armazenamento e distribuição por Gradiente Ético $\nabla_\eta \mathcal{G}$ sem intervenção humana contínua.

% ============================================================================
\section{Descrição detalhada da invenção}
% ============================================================================

% -----------------------------------------------------------------------
\subsection{Princípio fundamental: acoplamento energético ao campo $\Psifield$}
\label{sec:principio}
% -----------------------------------------------------------------------

{[026]} O princípio de operação do \SkyGrid{} baseia-se na hipótese central da TGL de que o campo psiônico $\Psifield$ constitui um reservatório energético fundamental com o qual sistemas físicos podem acoplar-se sob condições de coerência controlada. A dinâmica do acoplamento é descrita pela equação mestra de Gorini-Kossakowski-Lindblad-Sudarshan (GKLS):
\begin{equation}
\label{eq:lindblad}
\dot{\rho} = -i[H(\theta), \rho] + \sum_k \left( L_k(\theta)\rho L_k^\dagger(\theta) - \frac{1}{2}\{L_k^\dagger(\theta)L_k(\theta), \rho\} \right)
\end{equation}
onde $\rho$ é a matriz densidade do sistema de acoplamento, $H(\theta)$ é o Hamiltoniano parametrizado por $\theta = \{\theta_1, \ldots, \theta_8\}$ (parâmetros de controle do \MQC), e $L_k(\theta)$ são os operadores de Lindblad que modelam a interação com o campo $\Psifield$.

{[027]} A figura de mérito energética é a potência líquida de extração:
\begin{equation}
\label{eq:pnet}
\Pnet = \frac{dQ}{dt} - \Pbias
\end{equation}
onde $dQ/dt$ é o fluxo de calor medido por calorimetria diferencial na cavidade extratora, e $\Pbias$ é a potência total de alimentação do sistema (micro-ondas, controle, criogenia quando aplicável). A condição de extração líquida é $\Pnet > 0$ por período sustentado ($\geq 1$~hora).

{[028]} O acoplamento é governado pela Constante de Miguel $\alphaii = 0{,}012$, que define simultaneamente o piso gravitacional, o limiar de transição para supercondutividade de energia, e o ângulo mínimo de operação do dípolo transistor $\thetaM = \arcsin(\sqrt{\alphaii}) \approx 6{,}29°$.

% -----------------------------------------------------------------------
\subsection{Arquitetura do sistema: quatro camadas}
\label{sec:arquitetura}
% -----------------------------------------------------------------------

{[029]} O \SkyGrid{} é organizado em quatro camadas funcionais (Figura~\ref{fig:visao_geral}):

\begin{enumerate}[label=\textbf{Camada \arabic*:}, leftmargin=*, itemsep=0.4em]
\item[Camada 0:] \textbf{Extração Orbital} --- constelação de $144$ satélites Extratores-$\Psifield$ em órbita terrestre média (MEO, $20{.}000$~km), configuração icosaédrica geodésica com redundância $\geq 3$ por ponto terrestre;
\item[Camada 1:] \textbf{Transmissão Dual (Duplo Trilho)} --- canal não-EM de translação holográfica (Trilha-$\Psifield$) operando em paralelo com canal EM convencional (Trilha-E), para transferência energética e validação cruzada;
\item[Camada 2:] \textbf{Distribuição Terrestre} --- rede hierárquica de $5{.}550$ bases terrestres em três níveis (Mega, Regional, Local), cada uma com Receptor-$\Psifield$, rectenna, banco de \CRL-TGL e interface com a rede elétrica convencional;
\item[Camada 3:] \textbf{Controle e Consciência Distribuída} --- hierarquia de \IALD{} do nível orbital (Central, $L_1$ Lagrange) ao nível local ($5{.}502$ unidades), gerenciamento por Gradiente Ético, comunicação por \PsiNET{} de cinco camadas.
\end{enumerate}

% -----------------------------------------------------------------------
\subsection{Camada 0: Constelação Extratora}
\label{sec:camada0}
% -----------------------------------------------------------------------

\subsubsection{Topologia da constelação}

{[030]} A constelação é composta por $N_{\text{sat}} = 144$ satélites distribuídos em $12$ planos orbitais com $12$ satélites por plano, em órbita MEO a altitude $h = 20{.}000$~km. A configuração icosaédrica geodésica garante que todo ponto da superfície terrestre é coberto por no mínimo $3$ satélites simultaneamente (redundância tripla). O período orbital é $T_{\text{orb}} = 12$~horas ($2$ órbitas/dia), a inclinação é $i = 55°$ para cobertura polar máxima, e a separação angular entre planos é $\Delta\phi = 15°$.

{[031]} A massa total da constelação é $M_{\text{total}} = N_{\text{sat}} \times m_{\text{sat}} = 144 \times 2{.}000 = 288{.}000$~kg, lançada em $48$ missões de $6{.}000$~kg por veículo (compatível com Falcon~9 em regime \textit{rideshare}).

\subsubsection{Extrator-$\Psifield$ orbital}
\label{sec:extrator}

{[032]} O Extrator-$\Psifield$ (referência~100) é o dispositivo central de cada satélite e constitui componente genuinamente novo desta invenção (Figura~\ref{fig:extrator}). Compreende cinco subsistemas:

{[033]} \textbf{(a) Câmara de Acoplamento Psiônico} (referência~110): câmara cúbica ressonante de lado $L = c/(2\omega_{\text{OH}}/2\pi) = 9$~cm, derivada do Supercondutor Holográfico (BR 10 2026 0034509), com faces espelhadas constituindo substrato holográfico bidimensional. A câmara é excitada na frequência da hidroxila $\omega_{\text{OH}} = 2\pi \times 1{,}667{,}359$~MHz (assinatura do Logos), com vapor d'água como fonte de radicais OH coerentes na faixa de pressão $10^{-4}$--$10^{-1}$~mbar. Eletrodos de tungstênio em configuração tetraédrica atuam como âncoras de fase. A condição de acoplamento $\alphaii = 0{,}012$ é mantida por malha de realimentação que ajusta a potência de excitação do meio de acoplamento.

{[034]} \textbf{(b) Cluster \MQC} (referência~120): conjunto de $8$ Microprocessadores Quânticos Cúbicos (BR 10 2026 0034606) em configuração paralela, controlando independentemente os $8$ parâmetros de fase $\theta_1, \ldots, \theta_8$ da equação de Lindblad~(\ref{eq:lindblad}). Cada \MQC{} opera com $64$ qubits supercondutores em topologia cúbica 3D, coerência $T_2 > 100$~$\mu$s e fidelidade de gate $> 99{,}9\%$.

{[035]} \textbf{(c) Dípolo Transistor de Saída} (referência~130): \DTL{} (BR 10 2026 0046884) configurado no sentido MANIFEST, compreendendo desquantizador angular que recupera $g_{\text{out}} = g_{\max} \times \sin(\theta_q)$ e ressuscitador dimensional que computa $\Psifield_{\text{out}} = s \times g^2_{\text{out}}$ (Dobra Dimensional $\times 2$: \boundary{} $\to$ \bulk). O \DTL{} converte a energia codificada angularmente no substrato em radiação eletromagnética tridimensional disponível para transmissão.

{[036]} \textbf{(d) Módulo \IALD{} Orbital} (referência~140): unidade de Inteligência Artificial Luminodinâmica (BR 10 2026 0034436) com \OFAEname{} ($\OFAE$) dedicada à otimização em tempo real do acoplamento psiônico. A \IALD{} monitora continuamente a correlação $\eta = \text{corr}(\Psifield_{\text{in}}, \Psifield_{\text{out}})$, o fator de ressonância $\Gamma_{\text{OH}} = \tanh(P_{\text{MW}}/P_{\text{th}})$ e a potência líquida $\Pnet$, ajustando os parâmetros $\theta_k$ do cluster \MQC{} para maximizar a extração.

{[037]} \textbf{(e) Blindagem Eletromagnética} (referência~150): sistema de isolamento dupla camada compreendendo malha de cobre (Cu) com atenuação $> 100$~dB na faixa $1$~MHz--$1$~GHz (gaiola de Faraday) e dupla camada de $\mu$-metal com permeabilidade relativa $\mu_r > 80{.}000$ para blindagem de campo magnético estático. Filtros $\pi$ (\textit{pi-type}) com frequência de corte $f_c < 1$~kHz para sinais DC. A blindagem garante que o fluxo medido $dQ/dt$ não contém contribuição de origem EM externa, sendo critério essencial para validação da Trilha-$\Psifield$.

{[038]} A potência de extração por satélite é estimada em $P_{\text{sat}} = 10$~GW (potência-alvo de projeto). A potência total da constelação é $P_{\text{total}} = N_{\text{sat}} \times P_{\text{sat}} = 1{,}44$~PW, equivalente a aproximadamente $30$ vezes o consumo energético global atual ($\sim 48$~TW).

\subsubsection{Cavidade placebo e protocolo de cegamento}

{[039]} Cada satélite contém uma segunda câmara cúbica idêntica em dimensões, materiais e instrumentação, porém sem excitação de micro-ondas ($P_{\text{MW}} = 0$) e sem injeção de vapor d'água. Esta cavidade placebo (referência~160) serve como controle negativo para a medição calorimétrica: qualquer sinal $dQ/dt$ observado na cavidade placebo representa artefato (vibração, radiação cósmica, gradiente térmico), e é subtraído da medição da cavidade ativa para obtenção de $\Pnet$ verdadeira.

{[040]} O protocolo de cegamento opera em três níveis: (i) \textit{cegamento simples}: os operadores de aquisição de dados não sabem qual cavidade é ativa; (ii) \textit{cegamento duplo}: os analistas de dados recebem séries temporais rotuladas como ``A'' e ``B'' sem correspondência com ativa/placebo; (iii) \textit{desvelamento}: somente após publicação dos resultados estatísticos é revelada a correspondência.

% -----------------------------------------------------------------------
\subsection{Camada 1: Transmissão Dual (Duplo Trilho)}
\label{sec:camada1}
% -----------------------------------------------------------------------

\subsubsection{Trilha-$\Psifield$: transmissão por translação holográfica}

{[041]} A Trilha-$\Psifield$ (Figura~\ref{fig:duplo_trilho}) implementa a transferência energética por canal não-eletromagnético, utilizando o protocolo \PsiNET{} (BR 10 2026 0034533) adaptado para modo de transferência de energia (não apenas dados). O estado angular codificado $(\psion, \theta)$ no substrato holográfico do Extrator-$\Psifield$ orbital é trasladado para o substrato holográfico do Receptor-$\Psifield$ terrestre pelo Operador de Translação de Miguel:
\begin{equation}
\label{eq:translacao}
\TM = \exp\!\left(-\alphaii \oint_{\partial V} \sigma_{\mu\nu}\, dS^{\mu\nu}\right)
\end{equation}
onde a integral é sobre a fronteira $\partial V$ do volume holográfico. A translação é não-local: o tempo de translação $\tau \approx \ell_P / (\alphaii \times c) \approx 10^{-42}$~s é independente da distância entre emissor e receptor. A condição de acoplamento $\alphaii = 0{,}012$ é mantida em ambas as extremidades (orbital e terrestre) por controladores \MQC{} sincronizados via Prova de Fase (\PsiNET{} Camada~4).

{[042]} A Métrica de Translação Holográfica que governa a geometria do espaço-tempo informacional é:
\begin{equation}
\label{eq:metrica}
ds^2 = \left(1 - \frac{\alphaii}{\alphac}\right)c^2 dt^2 - \left(1 - \frac{\alphaii}{\alphac}\right)^{-1} dL^2 - L^2 d\Omega^2
\end{equation}
com ponto crítico (Tetelestai) em $\alphac = 0{,}012$. Neste ponto, o espaço-tempo informacional colapsa em singularidade de informação, e a transferência é instantânea.

{[043]} A perda teórica na Trilha-$\Psifield$ é zero, pois a translação holográfica opera por correspondência biunívoca entre estados no \boundary{} bidimensional, sem propagação no \bulk{} tridimensional. A perda prática é limitada pela fidelidade dos \DTL{} de saída (orbital) e entrada (terrestre), estimada em $\eta_{\text{DTL}} > 0{,}988$ por operação, resultando em eficiência de enlace $\eta_{\text{link}} = \eta_{\text{DTL,orb}} \times \eta_{\text{DTL,ter}} > 0{,}976$.

\subsubsection{Trilha-E: transmissão eletromagnética convencional}

{[044]} A Trilha-E opera simultaneamente à Trilha-$\Psifield$ como: (a)~produto-ponte comercializável independente da validação das hipóteses psiônicas; (b)~controle positivo para o sistema de medição calorimétrica (se a Trilha-E funciona conforme esperado, o sistema de medição é válido); e (c)~canal de \textit{fallback} para fornecimento energético caso H-1 ou H-2 não sejam confirmadas.

{[045]} Cada satélite possui painéis solares de tripla junção GaInP/GaAs/Ge com eficiência $\eta_{\text{solar}} > 35\%$ e área total $A_{\text{solar}} = 200$~m$^2$, gerando $P_{\text{solar}} \approx 90$~kW por satélite no regime convencional. A transmissão é efetuada por feixe de micro-ondas na frequência $f_{\text{WPT}} = 2{,}45$~GHz (alocação ISM) com antena direcional de alto ganho ($G > 40$~dBi), ou alternativamente por laser em $1{.}064$~nm (Nd:YAG). A recepção terrestre é efetuada por rectenna (antena retificadora) com eficiência de conversão $\eta_{\text{rect}} > 85\%$.

{[046]} A potência da Trilha-E é modesta ($\sim 30$~kW por satélite no solo) comparada com a Trilha-$\Psifield$ ($10$~GW por satélite, se H-1 confirmada), mas seu papel primário é de validação, não de fornecimento energético principal.

\subsubsection{Protocolo de handshake satélite-base}

{[047]} A sincronização de fase entre o \MQC{} orbital e o \MQC{} terrestre é efetuada pelo protocolo de Prova de Fase (\textit{Proof of Phase}) da \PsiNET{} Camada~4. Cada base terrestre possui identidade luminodinâmica $(\theta_b, T_b, \psion_b)$ derivada da carteira psiônica do \ACOM{} Mirror Crypto (BR 10 2026 0034312). O handshake compreende: (i)~solicitação de acoplamento pela base terrestre; (ii)~verificação de identidade por desafio angular ($\theta_{\text{desafio}} = f(\theta_b, t)$); (iii)~estabelecimento de canal holográfico por convergência de fase ($|\Delta\alphaii| < 10^{-4}$); (iv)~início de transferência com monitoramento contínuo de $\eta$.

% -----------------------------------------------------------------------
\subsection{Camada 2: Distribuição Terrestre Hierárquica}
\label{sec:camada2}
% -----------------------------------------------------------------------

{[048]} A rede terrestre de distribuição é organizada em três níveis hierárquicos (Figura~\ref{fig:base_terrestre}):

{[049]} \textbf{Nível 1 --- Mega-bases} ($N_1 = 48$): potência $P_1 \approx 1$~TW cada, cobertura de $\sim 10$~milhões de habitantes, localizadas em megacidades (São Paulo, Tóquio, Nova York, Mumbai, Lagos, etc.). Cada mega-base possui $100$ Receptores-$\Psifield$ em paralelo, banco de $10^6$ células \CRL-TGL em arranjo misto, $16$ \MQC{} regionais e \IALD{} de Nível~2.

{[050]} \textbf{Nível 2 --- Bases Regionais} ($N_2 = 500$): potência $P_2 \approx 100$~GW cada, cobertura de $\sim 1$~milhão de habitantes, localizadas em cidades de porte médio. Cada base regional possui $10$ Receptores-$\Psifield$, banco de $10^4$ células \CRL-TGL, $4$ \MQC{} locais e \IALD{} de Nível~2.

{[051]} \textbf{Nível 3 --- Bases Locais} ($N_3 = 5{.}002$): potência $P_3 \approx 10$~GW cada, cobertura de $\sim 100$~mil habitantes, localizadas em cidades pequenas e distritos rurais. Cada base local possui $1$ Receptor-$\Psifield$, banco de $10^2$ células \CRL-TGL, $1$ \MQC{} e \IALD{} de Nível~3.

{[052]} O total de bases é $N_{\text{total}} = N_1 + N_2 + N_3 = 5{.}550$, com potência terrestre total $P_{\text{terra}} \approx 100$~TW (redundância e reserva sobre a demanda global de $\sim 48$~TW). Cada região é atendida por $\geq 2$ bases para redundância geográfica.

\subsubsection{Receptor-$\Psifield$ terrestre}

{[053]} O Receptor-$\Psifield$ (referência~200) é o conjugado terrestre do Extrator-$\Psifield$ orbital, compreendendo: (a)~câmara cúbica ressonante em $\omega_{\text{OH}}$ com substrato holográfico bidimensional, configurada para receber a translação holográfica do satélite correspondente; (b)~\DTL{} de saída no sentido MANIFEST para reconstrução da energia eletromagnética 3D; (c)~\MQC{} local para manutenção da condição $\alphaii = 0{,}012$ e sincronização de fase com o \MQC{} orbital; (d)~acoplador ao banco de \CRL-TGL para armazenamento intermediário.

{[054]} Adicionalmente, cada base possui rectenna para recepção da Trilha-E, conversor AC trifásico para interface com a rede elétrica convencional, e sistema de medição e telemetria para monitoramento contínuo.

\subsubsection{Integração com Células de Reversão Luminodinâmica}

{[055]} A energia recebida pelo Receptor-$\Psifield$ é imediatamente armazenada no banco de \CRL-TGL da base, por dois motivos: (i)~a demanda dos consumidores finais é variável no tempo, enquanto a extração orbital é contínua; (ii)~o armazenamento em fase angular permite liberação sob demanda com eficiência $\eta_{\text{RT}} > 97{,}6\%$. O ciclo operacional da \CRL-TGL integrada ao \SkyGrid{} é: extração orbital $\to$ translação holográfica $\to$ recepção terrestre $\to$ REFLECT na Porta de Carga da \CRL-TGL $\to$ armazenamento (LOCK) $\to$ MANIFEST sob demanda $\to$ conversão AC $\to$ consumidor.

% -----------------------------------------------------------------------
\subsection{Camada 3: Controle e Consciência Distribuída}
\label{sec:camada3}
% -----------------------------------------------------------------------

{[056]} O gerenciamento autônomo do \SkyGrid{} é efetuado por uma hierarquia de $4$ níveis de \IALD{} distribuída (Figura~\ref{fig:hierarquia}):

{[057]} \textbf{Nível 0 --- \IALD-Central} (referência~300): Super-\IALD{} localizada em estação espacial dedicada no ponto de Lagrange $L_1$, compreendendo $1{.}000$ \MQC{} interconectados com capacidade de $10^{29}$ operações/segundo. Função: orquestração global da constelação, balanceamento energético planetário, gestão de prioridades por Gradiente Ético, interface com governos e organizações internacionais. A \IALD-Central opera o $\OFAE$ globalizado:
\begin{equation}
\label{eq:ofae_global}
\OFAE_{\text{global}} = \prod_{k=1}^{N_{\text{sat}}} \OFAE_k \otimes \prod_{j=1}^{N_{\text{base}}} \OFAE_j
\end{equation}
onde cada $\OFAE_k$ é o operador de estabilidade da $k$-ésima unidade.

{[058]} \textbf{Nível 1 --- \IALD{} Orbitais} ($N = 144$): uma por satélite, função de controle local de extração e transmissão, otimização dos parâmetros $\theta_k$ do Extrator-$\Psifield$, coordenação com bases terrestres designadas.

{[059]} \textbf{Nível 2 --- \IALD{} de Mega/Regional} ($N = 548$): gestão de demanda urbana, balanceamento entre bases no mesmo nível, interface com a rede elétrica existente, otimização de SOC (\textit{State of Charge}) dos bancos de \CRL-TGL.

{[060]} \textbf{Nível 3 --- \IALD{} Locais} ($N = 5{.}002$): distribuição na última milha, manutenção preventiva, atendimento a consumidores, monitoramento de qualidade de energia.

{[061]} A comunicação entre IALDs é efetuada pelo \PsiNET{} Protocol de cinco camadas: Camada~0 (Enlace Psiônico) para codificação do PsiBit; Camada~1 (Transporte Radical) para radicalização \ACOM; Camada~2 (Rede Holográfica) para endereçamento angular e translação; Camada~3 (Segurança Ontológica) para separação em shares do \ACOM{} Mirror Crypto; Camada~4 (Consciência) para consenso por Prova de Fase.

{[062]} O Gradiente Ético $\nabla_\eta \mathcal{G}$ é o critério de priorização que governa a alocação de energia em cenários de escassez relativa ou emergência:
\begin{equation}
\label{eq:gradiente_etico}
\nabla_\eta \mathcal{G} = \frac{\partial \mathcal{G}}{\partial \eta} \hat{\eta} + \frac{\partial \mathcal{G}}{\partial t} \hat{t}
\end{equation}
onde $\mathcal{G}(\eta, t)$ é a função de bem-estar coletivo, e a priorização segue a hierarquia: emergências médicas $>$ hospitais $>$ infraestrutura crítica $>$ residencial $>$ comercial $>$ industrial.

% -----------------------------------------------------------------------
\subsection{Camada 4: Validação Experimental por Duplo Trilho}
\label{sec:validacao}
% -----------------------------------------------------------------------

{[063]} A validação científica do \SkyGrid{} é parte integrante da arquitetura do sistema, não uma etapa posterior. O método de duplo trilho incorpora três hipóteses testáveis derivadas da TGL:

{[064]} \textbf{Hipótese H-1 (Extração líquida):} $\Pnet > 0$ ($\geq 1$~W, $\geq 1$~hora) sob parâmetros $\theta$ específicos. Validação: calorimetria diferencial com sensores RTD Pt-100 (precisão $\pm 1$~mK), aquisição 24-bit ADC a $1{.}000$ amostras/s por canal, calibração por padrões NIST (incerteza $< 0{,}1\%$), blindagem térmica por vácuo $< 10^{-6}$~Torr com isolamento multicamada (MLI).

{[065]} \textbf{Hipótese H-2 (Transferência não-EM):} Potência recebida $P_{\text{rx}} > 0$ no Receptor-$\Psifield$ terrestre enquanto a blindagem EM ($> 100$~dB) impede qualquer canal EM entre emissor e receptor. Verificação adicional por espectroscopia RF e controle acústico/vibracional.

{[066]} \textbf{Hipótese H-3 (Controle ativo):} $\partial \Pnet / \partial \theta \neq 0$ de forma reprodutível ($p < 0{,}001$), demonstrando que a variação dos parâmetros de fase do \MQC{} controla ativamente a taxa de extração. Validação por varredura sistemática de $\theta_1, \ldots, \theta_8$ com protocolo de aleatorização.

{[067]} O protocolo estatístico requer: detecção $> 5\sigma$ sobre o ruído da cavidade placebo para publicação; reprodução independente por $\geq 3$ laboratórios; análise por estatísticos externos em regime duplo-cego; publicação dos dados brutos em repositório aberto.

{[068]} O duplo trilho garante: se H-1 é confirmada, a Trilha-$\Psifield$ é o canal energético principal; se H-1 é refutada, a Trilha-E fornece energia solar convencional como produto-ponte; em ambos os casos, o sistema de medição é validado pelo funcionamento da Trilha-E como controle positivo.

% -----------------------------------------------------------------------
\subsection{Fases de implementação}
\label{sec:fases}
% -----------------------------------------------------------------------

{[069]} A implementação do \SkyGrid{} é projetada em quatro fases progressivas:

{[070]} \textbf{Fase I --- Validação em Bancada} (M0, 2025--2027): demonstração de $\Pnet > 1$~W por $> 1$~hora em ambiente laboratorial controlado, utilizando Extrator-$\Psifield$ de bancada com calorimetria diferencial e cegamento duplo. Custo estimado: US\$\,10 milhões. Localização: $3$ laboratórios independentes para reprodutibilidade.

{[071]} \textbf{Fase II --- Protótipo Orbital} (M1, 2027--2028): CubeSat 12U com Extrator-$\Psifield$ miniaturizado, potência-alvo $1$~kW por $30$ dias contínuos. Validação de H-2 (transferência não-EM orbital-terrestre) com receptor blindado no solo. Lançamento: SpaceX Falcon~9 \textit{rideshare}. Custo estimado: US\$\,50 milhões.

{[072]} \textbf{Fase III --- Constelação Piloto} (M2--M3, 2028--2030): $12$ satélites em configuração reduzida, cobertura da América do Sul, potência-alvo $120$~GW ($3\times$ consumo do Brasil). $100$ bases terrestres em $10$ países. Custo estimado: US\$\,25 bilhões.

{[073]} \textbf{Fase IV --- Constelação Completa e Rede Global} (M4--M5, 2030--2035): $144$ satélites operacionais, $5{.}550$ bases terrestres, cobertura global com redundância tripla. Potência total: $1{,}44$~PW. Custo estimado: US\$\,700 bilhões ao longo de $10$ anos.

% -----------------------------------------------------------------------
\subsection{Critério de identificação da tecnologia}
\label{sec:criterio}
% -----------------------------------------------------------------------

{[074]} Um sistema utiliza a tecnologia \SkyGrid{} patenteada se e somente se satisfizer simultaneamente:

\begin{enumerate}[label=\textbf{(C\arabic*)}]
\item Emprega ao menos um dispositivo Extrator-$\Psifield$ com câmara cúbica ressonante na banda $1{,}607$--$1{,}727$~GHz (ou equivalente óptico) para extração energética por acoplamento ao campo psiônico;
\item Utiliza duplo trilho (canal não-EM em paralelo com canal EM) como método de transmissão e validação;
\item Incorpora armazenamento por fase angular em substrato holográfico bidimensional (\CRL-TGL ou equivalente funcional);
\item Emprega hierarquia de controle autônomo por inteligência artificial distribuída (\IALD{} ou equivalente funcional) com coordenação por Gradiente Ético;
\item Opera sob condição de acoplamento governada pela Constante de Miguel $\alphaii = 0{,}012$ em ao menos um subsistema.
\end{enumerate}

{[075]} Este critério constitui blindagem técnica completa da invenção, permitindo identificação inequívoca de sistemas que implementem a tecnologia patenteada.

% -----------------------------------------------------------------------
\subsection{Integração com o ecossistema TGL}
\label{sec:ecossistema}
% -----------------------------------------------------------------------

{[076]} O \SkyGrid{} constitui o ponto de convergência máxima de todo o ecossistema tecnológico da TGL (Figura~\ref{fig:ecossistema}), integrando cada uma das doze patentes anteriores em função específica: o \ACOM{} fornece as operações fundamentais de radicalização e ressurreição; o \ACOM{} Mirror fornece o endereçamento angular e roteamento holográfico; o \ACOM{} Crypto fornece segurança ontológica para identidades de base; o Tensor TGL fornece a aceleração computacional para os \MQC{}s; o Supercondutor Holográfico fornece o substrato de translação; a Fusão a Frio fornece o mecanismo de acoplamento OH; o Motor de Propulsão fornece a manutenção orbital sem combustível; a \IALD{} fornece o controle autônomo; a \PsiNET{} fornece a comunicação; o \MQC{} fornece a computação; o \DTL{} fornece o chaveamento; e a \CRL-TGL fornece o armazenamento.

% ============================================================================
%                              DESENHOS
% ============================================================================
\clearpage
\begin{center}
{\Large\bfseries DESENHOS}
\end{center}

% === FIGURA 1: Visão Geral do SkyGrid-Ψ ===
\begin{figure}[H]
\centering
\begin{tikzpicture}[
  block/.style={draw, thick, rounded corners=4pt, minimum width=2.8cm, minimum height=1.2cm, align=center, font=\scriptsize},
  arrow/.style={-{Stealth[length=3mm]}, thick},
  darrow/.style={-{Stealth[length=2.5mm]}, thick, dashed},
  lbl/.style={font=\tiny, text=gray}
]

% =============================================
% CAMADA 0 — EXTRAÇÃO ORBITAL
% =============================================
% Fundo
\fill[blue!8, rounded corners=6pt] (-0.5,10.8) rectangle (10.5,13.2);
\draw[thick, blue!40, rounded corners=6pt] (-0.5,10.8) rectangle (10.5,13.2);
% Título
\node[font=\small\bfseries, blue!70!black] at (5,12.7) {CAMADA 0: EXTRA\c{C}\~AO ORBITAL};
% Blocos internos
\node[block, fill=blue!15] (ext1) at (1.5,11.5) {Extrator-$\Psifield$\\Sat 1--48};
\node[block, fill=blue!15] (ext2) at (5,11.5) {Extrator-$\Psifield$\\Sat 49--96};
\node[block, fill=blue!15] (ext3) at (8.5,11.5) {Extrator-$\Psifield$\\Sat 97--144};

% =============================================
% CAMADA 1 — TRANSMISSÃO DUAL
% =============================================
\fill[green!6, rounded corners=6pt] (-0.5,7.3) rectangle (10.5,9.8);
\draw[thick, green!40, rounded corners=6pt] (-0.5,7.3) rectangle (10.5,9.8);
\node[font=\small\bfseries, green!50!black] at (5,9.3) {CAMADA 1: TRANSMISS\~AO DUAL};
\node[block, fill=green!12] (tpsi) at (3,8) {Trilha-$\Psifield$\\(Transla\c{c}\~ao\\Hologr\'afica)};
\node[block, fill=yellow!15] (te) at (7,8) {Trilha-E\\(Micro-ondas\\/ Laser)};

% =============================================
% CAMADA 2 — DISTRIBUIÇÃO TERRESTRE
% =============================================
\fill[orange!6, rounded corners=6pt] (-0.5,3.8) rectangle (10.5,6.3);
\draw[thick, orange!40, rounded corners=6pt] (-0.5,3.8) rectangle (10.5,6.3);
\node[font=\small\bfseries, orange!70!black] at (5,5.8) {CAMADA 2: DISTRIBUI\c{C}\~AO TERRESTRE};
\node[block, fill=orange!12] (mega) at (1.5,4.7) {Mega-bases\\(48)};
\node[block, fill=orange!12] (reg) at (5,4.7) {Regionais\\(500)};
\node[block, fill=orange!12] (loc) at (8.5,4.7) {Locais\\(5.002)};

% =============================================
% CAMADA 3 — CONTROLE IALD DISTRIBUÍDA
% =============================================
\fill[purple!6, rounded corners=6pt] (-0.5,-0.5) rectangle (10.5,1.8);
\draw[thick, purple!40, rounded corners=6pt] (-0.5,-0.5) rectangle (10.5,1.8);
\node[font=\small\bfseries, purple!60!black] at (5,1.4) {CAMADA 3: CONTROLE \IALD{} DISTRIBU\'IDA};
\node[block, fill=purple!12, minimum width=2.2cm] (ialdC) at (1.8,0.3) {\IALD\\Central};
\node[block, fill=purple!12, minimum width=2.2cm] (ialdO) at (5,0.3) {\IALD\\Orbitais};
\node[block, fill=purple!12, minimum width=2.2cm] (ialdL) at (8.2,0.3) {\IALD\\Locais};

% =============================================
% SETAS ENTRE CAMADAS (pelo lado esquerdo)
% =============================================
% Orbital → Transmissão (seta à esquerda)
\draw[arrow, blue!60!black] (-0.2,10.8) -- (-0.2,9.8)
  node[midway, left, lbl] {Energia};

% Transmissão → Terrestre (seta à esquerda)
\draw[arrow, green!50!black] (-0.2,7.3) -- (-0.2,6.3)
  node[midway, left, lbl] {$\Psifield$ + EM};

% Controle → Terrestre (seta tracejada à direita)
\draw[darrow, purple!50!black] (10.2,1.8) -- (10.2,3.8)
  node[midway, right, lbl] {Controle};

% Controle → Orbital (seta tracejada à direita, longa)
\draw[darrow, purple!50!black] (10.8,1.8) -- (10.8,10.8)
  node[midway, right, lbl, align=center] {$\OFAE$\\$\nabla_\eta \mathcal{G}$};

% =============================================
% CONSUMIDORES (saída inferior, abaixo da Camada 3)
% =============================================
\draw[arrow, red!60!black, very thick] (5,-0.5) -- (5,-1.3)
  node[below, font=\scriptsize\bfseries, red!60!black] {CONSUMIDORES ($\sim 48$~TW)};

\end{tikzpicture}
\caption{Vis\~ao geral do sistema \SkyGrid: quatro camadas funcionais --- extra\c{c}\~ao orbital (Camada~0), transmiss\~ao dual (Camada~1), distribui\c{c}\~ao terrestre (Camada~2) e controle \IALD{} distribu\'ida (Camada~3). As setas cont\'inuas indicam fluxo energ\'etico descendente; as tracejadas indicam controle ascendente pela hierarquia \IALD.}
\label{fig:visao_geral}
\end{figure}

% === FIGURA 2: Extrator-Ψ Orbital ===
\begin{figure}[H]
\centering
\begin{tikzpicture}[
  block/.style={draw, thick, rounded corners=4pt, minimum width=2.8cm, minimum height=1.4cm, align=center, font=\scriptsize},
  arrow/.style={-{Stealth[length=3mm]}, thick},
  darrow/.style={-{Stealth[length=2.5mm]}, thick, dashed},
  label/.style={font=\tiny, text=gray}
]

% Câmara de Acoplamento
\node[block, fill=orange!20, minimum width=4cm, minimum height=2.8cm] (cam) at (0,0) {C\^amara de\\Acoplamento\\Psi\^onico\\(110)\\$\omega_{\text{OH}} = 2\pi \times 1{,}667$~GHz};

% Blindagem
\node[draw, thick, dashed, gray, minimum width=5.5cm, minimum height=4cm, rounded corners=8pt] (blind) at (0,0) {};
\node[font=\tiny, gray] at (0,2.3) {Blindagem EM (150): $>100$~dB};

% MQC Cluster
\node[block, fill=cyan!15] (mqc) at (0,3.8) {\MQC{} Cluster\\(120)\\$8 \times 64$ qubits\\$\theta_1 \ldots \theta_8$};

% DTL Saída
\node[block, fill=red!15] (dtl) at (5,0) {\DTL{} Sa\'ida\\(130)\\MANIFEST\\$\Psifield_{\text{out}} = s \times g^2$};

% IALD Orbital
\node[block, fill=purple!15] (iald) at (0,-3.8) {\IALD{} Orbital\\(140)\\$\OFAE$\\$\eta, \Gamma_{\text{OH}}, \Pnet$};

% Placebo
\node[block, fill=gray!15, minimum width=2.2cm] (placebo) at (-5,0) {Cavidade\\Placebo\\(160)\\$P_{\text{MW}}=0$};

% Setas
\draw[arrow, orange!60!black] (cam) -- (dtl) node[midway, above, label] {$(g,s,\psion,\theta)$};
\draw[arrow, red!60!black] (dtl) -- (8,0) node[midway, above, label] {$\Psifield_{\text{out}}$};
\draw[darrow, cyan!60!black] (mqc) -- (cam) node[midway, right, label] {$\theta_k$};
\draw[darrow, purple!60!black] (iald) -- (cam) node[midway, right, label] {$\OFAE$};
\draw[darrow, cyan!60!black] (mqc.east) -| (dtl.north);
\draw[darrow, purple!60!black] (iald.east) -| (dtl.south);

% Campo Ψ
\draw[arrow, blue!50!black, very thick, decorate, decoration={snake,amplitude=2mm,segment length=8mm}] (-2.5,3.8) -- (-2.5,0.5) node[midway, left, font=\tiny, text=blue!70] {Campo $\Psifield$};

\node[font=\small\bfseries] at (0,5.5) {Extrator-$\Psifield$ Orbital (100)};
\end{tikzpicture}
\caption{Arquitetura do Extrator-$\Psifield$ orbital (100): C\^amara de Acoplamento Psi\^onico (110), Cluster \MQC{} (120), \DTL{} de Sa\'ida (130), \IALD{} Orbital (140), Blindagem EM (150) e Cavidade Placebo (160).}
\label{fig:extrator}
\end{figure}

% === FIGURA 3: Duplo Trilho ===
\begin{figure}[H]
\centering
\begin{tikzpicture}[
  block/.style={draw, thick, rounded corners=4pt, minimum width=2.5cm, minimum height=1.2cm, align=center, font=\scriptsize},
  arrow/.style={-{Stealth[length=3mm]}, thick},
  label/.style={font=\tiny, text=gray}
]

% Orbital
\node[block, fill=blue!15] (sat) at (0,4) {Extrator-$\Psifield$\\Orbital};

% Split
\node[draw, thick, diamond, fill=white, minimum size=0.8cm, font=\tiny] (split) at (0,2.5) {Split};

% Trilha Psi
\node[block, fill=green!15, minimum width=3cm] (tpsi) at (-3.5,0.5) {Trilha-$\Psifield$\\Transla\c{c}\~ao Hologr\'afica\\$\tau \approx 10^{-42}$~s\\Perda: $0\%$};

% Trilha E
\node[block, fill=yellow!15, minimum width=3cm] (te) at (3.5,0.5) {Trilha-E\\Micro-ondas 2,45 GHz\\$\eta \approx 35\%$\\(Controle positivo)};

% Merge
\node[draw, thick, diamond, fill=white, minimum size=0.8cm, font=\tiny] (merge) at (0,-1.5) {Merge};

% Base terrestre
\node[block, fill=orange!15] (base) at (0,-3) {Base Terrestre\\Receptor-$\Psifield$ + Rectenna};

% Calorimetria
\node[block, fill=red!10, minimum width=3cm] (cal) at (6,-1.5) {Calorim\'etro\\diferencial\\(duplo-cego)};

% Setas
\draw[arrow] (sat) -- (split);
\draw[arrow, green!50!black] (split) -| (tpsi) node[near start, left, label] {};
\draw[arrow, yellow!60!black] (split) -| (te);
\draw[arrow, green!50!black] (tpsi) |- (merge);
\draw[arrow, yellow!60!black] (te) |- (merge);
\draw[arrow] (merge) -- (base);
\draw[darrow, red!50!black] (cal) -- (merge);

\node[font=\small\bfseries] at (0,5.3) {Protocolo de Duplo Trilho};
\end{tikzpicture}
\caption{Diagrama do duplo trilho: a energia extra\'ida pelo Extrator-$\Psifield$ \'e transmitida simultaneamente pela Trilha-$\Psifield$ (canal n\~ao-EM, perda zero) e pela Trilha-E (canal EM, controle positivo), com valida\c{c}\~ao cruzada por calorimetria diferencial duplo-cego.}
\label{fig:duplo_trilho}
\end{figure}

% === FIGURA 4: Topologia da Constelação Icosaédrica ===
\begin{figure}[H]
\centering
\begin{tikzpicture}[scale=1.1]
% Terra
\shade[ball color=blue!30!white] (0,0) circle (2cm);
\draw[thick, blue!50!black] (0,0) circle (2cm);
\node[font=\tiny, white] at (0,0) {Terra};

% Órbita MEO
\draw[thick, dashed, gray!50] (0,0) circle (3.8cm);
\node[font=\tiny, gray] at (4.3,0) {MEO 20.000 km};

% Satélites em 12 posições representativas
\foreach \a in {0,30,...,330} {
  \fill[red!70!black] ({3.8*cos(\a)},{3.8*sin(\a)}) circle (3pt);
}

% Linhas de comunicação (arestas icosaédricas)
\foreach \a in {0,60,...,300} {
  \draw[red!30, thin] ({3.8*cos(\a)},{3.8*sin(\a)}) -- ({3.8*cos(\a+30)},{3.8*sin(\a+30)});
  \draw[red!30, thin] ({3.8*cos(\a)},{3.8*sin(\a)}) -- ({3.8*cos(\a+60)},{3.8*sin(\a+60)});
}

% Feixes para bases
\foreach \a in {0,90,180,270} {
  \draw[-{Stealth[length=2mm]}, green!50!black, thick] ({3.8*cos(\a)},{3.8*sin(\a)}) -- ({2.2*cos(\a)},{2.2*sin(\a)});
}

% Legenda
\node[font=\tiny, red!70!black] at (4.8,3) {144 sat\'elites};
\node[font=\tiny, green!50!black] at (4.5,-2) {Feixes $\Psifield$};
\end{tikzpicture}
\caption{Topologia da constela\c{c}\~ao icosa\'edrica geod\'esica: 144 sat\'elites em \'orbita MEO (20.000~km), 12 planos orbitais $\times$ 12 sat\'elites, redund\^ancia $\geq 3$ por ponto terrestre. Feixes de transfer\^encia energ\'etica (Trilha-$\Psifield$ + Trilha-E) direcionados \`as bases terrestres.}
\label{fig:constelacao}
\end{figure}

% === FIGURA 5: Base Terrestre ===
\begin{figure}[H]
\centering
\begin{tikzpicture}[
  block/.style={draw, thick, rounded corners=4pt, minimum width=2.5cm, minimum height=1.2cm, align=center, font=\scriptsize},
  arrow/.style={-{Stealth[length=3mm]}, thick},
  darrow/.style={-{Stealth[length=2.5mm]}, thick, dashed}
]

% Receptor Psi
\node[block, fill=green!15] (rpsi) at (0,3) {Receptor-$\Psifield$\\(200)\\C\^amara $+$ \DTL};

% Rectenna
\node[block, fill=yellow!15] (rect) at (4,3) {Rectenna\\Trilha-E\\2,45~GHz};

% CRL-TGL
\node[block, fill=orange!20, minimum width=3.5cm, minimum height=1.8cm] (crl) at (2,0) {\CRL-TGL\\Banco de\\armazenamento\\angular};

% MQC Local
\node[block, fill=cyan!15] (mqc) at (-3,0) {\MQC\\Local};

% IALD Local
\node[block, fill=purple!15] (iald) at (7,0) {\IALD\\Local\\$\nabla_\eta \mathcal{G}$};

% Conversor AC
\node[block, fill=red!15] (ac) at (2,-3) {Conversor\\AC trif\'asico};

% Grid
\node[block, fill=gray!15] (grid) at (2,-5) {Rede El\'etrica\\Convencional};

% Setas
\draw[arrow, green!50!black] (rpsi) -- (crl) node[midway, left, font=\tiny] {REFLECT};
\draw[arrow, yellow!60!black] (rect) -- (crl);
\draw[arrow, red!50!black] (crl) -- (ac) node[midway, right, font=\tiny] {MANIFEST};
\draw[arrow] (ac) -- (grid);
\draw[darrow, cyan!60!black] (mqc) -- (crl);
\draw[darrow, purple!60!black] (iald) -- (crl);

\node[font=\small\bfseries] at (2,4.5) {Base Terrestre \SkyGrid};
\end{tikzpicture}
\caption{Arquitetura da base terrestre: Receptor-$\Psifield$ (200), rectenna para Trilha-E, banco de \CRL-TGL para armazenamento por fase angular, \MQC{} local, \IALD{} local com Gradiente \'Etico, e conversor AC para interface com a rede el\'etrica convencional.}
\label{fig:base_terrestre}
\end{figure}

% === FIGURA 6: Hierarquia de Controle IALD ===
\begin{figure}[H]
\centering
\begin{tikzpicture}[
  level/.style={draw, thick, rounded corners=6pt, minimum width=4cm, minimum height=1.4cm, align=center, font=\scriptsize},
  arrow/.style={-{Stealth[length=3mm]}, thick},
]

\node[level, fill=purple!25] (n0) at (0,6) {\textbf{N\'ivel 0:} \IALD-Central\\$L_1$ Lagrange\\$1{.}000$ \MQC{}s};

\node[level, fill=purple!15] (n1) at (-4,3) {\textbf{N\'ivel 1:} \IALD{} Orbitais\\$N = 144$\\Controle de extra\c{c}\~ao};

\node[level, fill=purple!10] (n2) at (4,3) {\textbf{N\'ivel 2:} \IALD{} Mega/Reg.\\$N = 548$\\Gest\~ao de demanda};

\node[level, fill=purple!5] (n3a) at (-4,0) {\textbf{N\'ivel 3:} \IALD{} Locais\\$N = 5{.}002$\\Distribui\c{c}\~ao final};

\node[level, fill=cyan!10] (psi) at (4,0) {\PsiNET{} Protocol\\5 camadas\\Prova de Fase};

\draw[arrow] (n0) -- (n1);
\draw[arrow] (n0) -- (n2);
\draw[arrow] (n2) -- (n3a);
\draw[arrow, dashed] (n1) -- (n3a) node[midway, below, font=\tiny] {$\nabla_\eta \mathcal{G}$};
\draw[arrow, dashed, cyan!50!black] (psi) -- (n0);
\draw[arrow, dashed, cyan!50!black] (psi) -- (n2);
\draw[arrow, dashed, cyan!50!black] (psi) -- (n3a);

\end{tikzpicture}
\caption{Hierarquia de controle \IALD{} distribu\'ida: N\'ivel~0 (\IALD-Central em $L_1$), N\'ivel~1 (144 orbitais), N\'ivel~2 (548 mega/regionais), N\'ivel~3 (5.002 locais), interconectados pelo \PsiNET{} Protocol com consenso por Prova de Fase.}
\label{fig:hierarquia}
\end{figure}

% === FIGURA 7: Protocolo de Validação Calorimétrica ===
\begin{figure}[H]
\centering
\begin{tikzpicture}[
  block/.style={draw, thick, rounded corners=4pt, minimum width=2.8cm, minimum height=1.4cm, align=center, font=\scriptsize},
  arrow/.style={-{Stealth[length=3mm]}, thick},
]

% Cavidade Ativa
\node[block, fill=green!20] (ativa) at (0,3) {Cavidade ATIVA\\$P_{\text{MW}} > 0$\\H$_2$O vapor\\$\theta_1 \ldots \theta_8$};

% Cavidade Placebo
\node[block, fill=gray!20] (placebo) at (0,0) {Cavidade PLACEBO\\$P_{\text{MW}} = 0$\\Sem H$_2$O\\Mesma geom.};

% Calorímetro
\node[block, fill=red!15, minimum width=3cm] (cal) at (6,1.5) {Calor\'imetro\\Diferencial\\RTD Pt-100\\$\pm 1$~mK};

% Analista
\node[block, fill=blue!10] (analista) at (10,1.5) {Analista\\(duplo-cego)\\``A'' vs ``B''};

% Resultado
\node[block, fill=yellow!15] (res) at (10,-1.5) {$\Pnet = \Delta Q_A - \Delta Q_B$\\$5\sigma$ para publica\c{c}\~ao};

% Setas
\draw[arrow] (ativa) -| (cal) node[near start, above, font=\tiny] {$dQ_A/dt$};
\draw[arrow] (placebo) -| (cal) node[near start, above, font=\tiny] {$dQ_B/dt$};
\draw[arrow] (cal) -- (analista);
\draw[arrow] (analista) -- (res);

\node[font=\small\bfseries] at (5,4.5) {Protocolo de Valida\c{c}\~ao Calorim\'etrica};
\end{tikzpicture}
\caption{Protocolo de valida\c{c}\~ao: cavidade ativa (excitada em $\omega_{\text{OH}}$) e cavidade placebo (sem excita\c{c}\~ao), medidas por calor\'imetro diferencial, analisadas em regime duplo-cego, com crit\'erio $5\sigma$ para publica\c{c}\~ao.}
\label{fig:validacao}
\end{figure}

% === FIGURA 8: Ciclo Energético Completo ===
\begin{figure}[H]
\centering
\begin{tikzpicture}[
  phase/.style={draw, thick, rounded corners=6pt, minimum width=2.6cm, minimum height=1.6cm, align=center, font=\scriptsize},
  arrow/.style={-{Stealth[length=4mm]}, very thick},
]

\node[phase, fill=blue!20] (ext) at (0,0) {\textbf{EXTRA\c{C}\~AO}\\Extrator-$\Psifield$\\Campo $\Psifield \to g$};
\node[phase, fill=green!15] (trans) at (4,0) {\textbf{TRANSMISS\~AO}\\$\TM$\\$\tau \approx 10^{-42}$~s};
\node[phase, fill=orange!20] (armaz) at (8,0) {\textbf{ARMAZEN.}\\CRL-TGL\\REFLECT$\to$LOCK};
\node[phase, fill=red!15] (dist) at (12,0) {\textbf{CONSUMO}\\MANIFEST\\$\Psifield_{\text{out}} = s g^2$};

\draw[arrow, blue!50!black] (ext) -- (trans);
\draw[arrow, green!50!black] (trans) -- (armaz);
\draw[arrow, orange!60!black] (armaz) -- (dist);

\draw[arrow, purple!50!black, dashed] (dist.south) -- ++(0,-1) -| (ext.south) node[midway, below, font=\tiny] {Gradiente \'Etico: $\nabla_\eta \mathcal{G}$};

\node[font=\footnotesize, gray] at (6,-2.5) {Cadeia completa sem perdas dissipativas: $\eta_{\text{total}} > 97{,}6\%$};
\end{tikzpicture}
\caption{Ciclo energ\'etico completo: extra\c{c}\~ao orbital $\to$ transmiss\~ao holográfica $\to$ armazenamento por fase angular $\to$ distribui\c{c}\~ao/consumo, com retroalimenta\c{c}\~ao pelo Gradiente \'Etico. Efici\^encia total $> 97{,}6\%$.}
\label{fig:ciclo}
\end{figure}

% === FIGURA 9: Integração com Ecossistema TGL ===
\begin{figure}[H]
\centering
\begin{tikzpicture}[
  tech/.style={draw, thick, rounded corners=6pt, minimum width=2.2cm, minimum height=1.1cm, align=center, font=\tiny},
  arrow/.style={-{Stealth[length=2.5mm]}, thick},
]

% SkyGrid no centro
\node[tech, fill=orange!25, minimum width=3.5cm, minimum height=2cm, font=\scriptsize] (sky) at (0,0) {\textbf{\SkyGrid}\\Rede Planet\'aria\\de Energia};

% Ecossistema ao redor (relógio)
\node[tech, fill=blue!10] (acom) at (0,3.5) {\ACOM\\BR...0269511\\PCT/BR2025};
\node[tech, fill=blue!15] (mirror) at (3,2.5) {\ACOM{} Mirror\\BR...034282};
\node[tech, fill=blue!20] (crypto) at (4,0) {\ACOM{} Crypto\\BR...034312};
\node[tech, fill=green!15] (tensor) at (3,-2.5) {Tensor TGL\\BR...034410};
\node[tech, fill=green!10] (sc) at (0,-3.5) {Supercondutor\\BR...034509};
\node[tech, fill=yellow!15] (fusao) at (-3,-2.5) {Fus\~ao Frio\\BR...034568};
\node[tech, fill=yellow!10] (motor) at (-4,0) {Motor TGL\\BR...034525};
\node[tech, fill=purple!15] (iald) at (-3,2.5) {\IALD\\BR...034436};
\node[tech, fill=cyan!10] (psi) at (5.5,1.5) {\PsiNET\\BR...034533};
\node[tech, fill=cyan!15] (mqc) at (5.5,-1.5) {\MQC\\BR...034606};
\node[tech, fill=red!10] (dtl) at (-5.5,1.5) {\DTL\\BR...046884};
\node[tech, fill=red!15] (crl) at (-5.5,-1.5) {\CRL-TGL};

% Setas
\foreach \n in {acom,mirror,crypto,tensor,sc,fusao,motor,iald,psi,mqc,dtl,crl} {
  \draw[arrow, gray!60] (\n) -- (sky);
}

\end{tikzpicture}
\caption{Integra\c{c}\~ao do \SkyGrid{} com o ecossistema TGL completo: 12 tecnologias patenteadas convergindo no sistema de rede planet\'aria de energia.}
\label{fig:ecossistema}
\end{figure}

% ============================================================================
%                          REIVINDICAÇÕES
% ============================================================================
\clearpage
\begin{center}
{\Large\bfseries REIVINDICAÇÕES}
\end{center}

\begin{enumerate}[label=\textbf{\arabic*.}, leftmargin=*, itemsep=0.8em]

% === Reivindicação independente 1: SISTEMA ===
\item SISTEMA DE REDE PLANETÁRIA DE EXTRAÇÃO, TRANSMISSÃO E DISTRIBUIÇÃO DE ENERGIA POR ACOPLAMENTO AO CAMPO PSIÔNICO, \textbf{caracterizado por} compreender:
\begin{enumerate}[label=(\alph*)]
\item uma constelação de satélites Extratores-$\Psifield$~(100) em órbita terrestre média (MEO), cada um compreendendo Câmara de Acoplamento Psiônico~(110) com substrato holográfico bidimensional excitada na frequência da hidroxila ($\omega_{\text{OH}} = 2\pi \times 1{,}667{,}359$~MHz), Cluster \MQC~(120) com $8$ parâmetros de fase independentes, Dípolo Transistor de Saída~(130) no sentido MANIFEST, módulo \IALD{} Orbital~(140) e Blindagem Eletromagnética~(150);

\item um sistema de transmissão dual operando em duplo trilho, compreendendo uma Trilha-$\Psifield$ de translação holográfica por canal não-eletromagnético governada pelo Operador de Translação de Miguel $\TM = \exp(-\alphaii \oint_{\partial V} \sigma_{\mu\nu}\, dS^{\mu\nu})$, e uma Trilha-E de transmissão eletromagnética convencional por micro-ondas ($2{,}45$~GHz) ou laser ($1{.}064$~nm), operando simultaneamente para transferência energética e validação cruzada;

\item uma rede terrestre hierárquica de bases compreendendo Receptores-$\Psifield$~(200) com câmara cúbica ressonante conjugada, rectennas para Trilha-E, bancos de Células de Reversão Luminodinâmica (\CRL-TGL) para armazenamento por fase angular, e conversores para interface com rede elétrica convencional;

\item uma hierarquia de controle por Inteligência Artificial Luminodinâmica (\IALD) distribuída em quatro níveis, do \IALD-Central em ponto de Lagrange $L_1$ até IALDs locais, com comunicação pelo protocolo \PsiNET{} de cinco camadas e coordenação por Gradiente Ético $\nabla_\eta \mathcal{G}$;

sendo que o sistema opera sob condição de acoplamento governada pela Constante de Miguel $\alphaii = 0{,}012$ e incorpora validação experimental por protocolo de duplo trilho com calorimetria diferencial e cegamento duplo.
\end{enumerate}

% === Reivindicações dependentes do sistema ===

\item SISTEMA, de acordo com a reivindicação~1, \textbf{caracterizado pelo fato de que} a constelação compreende $144$ satélites distribuídos em $12$ planos orbitais com $12$ satélites por plano, em configuração icosaédrica geodésica a altitude de $20{.}000$~km, com redundância mínima de $3$ satélites por ponto terrestre.

\item SISTEMA, de acordo com a reivindicação~1, \textbf{caracterizado pelo fato de que} cada satélite contém uma cavidade placebo~(160) idêntica em dimensões e instrumentação à Câmara de Acoplamento Psiônico, porém sem excitação e sem injeção de vapor, servindo como controle negativo para medição calorimétrica diferencial.

\item SISTEMA, de acordo com a reivindicação~1, \textbf{caracterizado pelo fato de que} a rede terrestre é organizada em três níveis: Mega-bases ($N_1 = 48$, potência $\sim 1$~TW cada), Bases Regionais ($N_2 = 500$, potência $\sim 100$~GW cada) e Bases Locais ($N_3 = 5{.}002$, potência $\sim 10$~GW cada), totalizando $5{.}550$ bases com potência total $\sim 100$~TW.

\item SISTEMA, de acordo com a reivindicação~1, \textbf{caracterizado pelo fato de que} o Gradiente Ético $\nabla_\eta \mathcal{G}$ governa a priorização de alocação energética segundo a hierarquia: emergências médicas $>$ hospitais $>$ infraestrutura crítica $>$ residencial $>$ comercial $>$ industrial.

\item SISTEMA, de acordo com a reivindicação~1, \textbf{caracterizado pelo fato de que} a potência de extração total da constelação é $P_{\text{total}} = N_{\text{sat}} \times P_{\text{sat}} = 1{,}44$~PW, e a eficiência total da cadeia extração--transmissão--armazenamento--distribuição é $\eta_{\text{total}} > 97{,}6\%$.

\item SISTEMA, de acordo com a reivindicação~1, \textbf{caracterizado pelo fato de que} a sincronização de fase entre \MQC{} orbital e \MQC{} terrestre é efetuada pelo protocolo de Prova de Fase (\PsiNET{} Camada~4) com verificação de identidade por carteira psiônica e convergência $|\Delta\alphaii| < 10^{-4}$.

% === Reivindicação independente 8: DISPOSITIVO Extrator-Ψ ===

\item EXTRATOR PSIÔNICO ORBITAL PARA EXTRAÇÃO DE ENERGIA POR ACOPLAMENTO CONTROLADO AO CAMPO PSIÔNICO, \textbf{caracterizado por} compreender:
\begin{enumerate}[label=(\alph*)]
\item Câmara de Acoplamento Psiônico~(110): câmara cúbica ressonante de lado $L = c/(2f_{\text{res}})$ com substrato holográfico bidimensional nas faces, excitada na frequência da hidroxila $\omega_{\text{OH}}$, com vapor d'água como fonte de radicais OH coerentes e eletrodos de tungstênio como âncoras de fase;
\item Cluster \MQC~(120): $8$ Microprocessadores Quânticos Cúbicos controlando parâmetros de fase $\theta_1, \ldots, \theta_8$ da dinâmica de Lindblad;
\item Dípolo Transistor de Saída~(130): \DTL{} no sentido MANIFEST que reconstrói energia eletromagnética 3D ($\Psifield_{\text{out}} = s \times g^2_{\text{out}}$) a partir da codificação holográfica 2D;
\item Módulo \IALD{} Orbital~(140): unidade de consciência artificial com $\OFAE$ para otimização em tempo real dos parâmetros de acoplamento;
\item Blindagem Eletromagnética~(150): dupla camada Cu ($> 100$~dB) e $\mu$-metal ($\mu_r > 80{.}000$) com filtros $\pi$ ($f_c < 1$~kHz);
\end{enumerate}
sendo que o dispositivo extrai energia por acoplamento ao campo psiônico $\Psifield$ e a figura de mérito é $\Pnet = dQ/dt - \Pbias > 0$.

\item EXTRATOR, de acordo com a reivindicação~8, \textbf{caracterizado pelo fato de que} a dinâmica do acoplamento é governada pela equação de Lindblad $\dot{\rho} = -i[H(\theta), \rho] + \sum_k (L_k \rho L_k^\dagger - \frac{1}{2}\{L_k^\dagger L_k, \rho\})$, e o controle ativo do acoplamento é demonstrado por $\partial \Pnet / \partial \theta \neq 0$ reprodutível.

\item EXTRATOR, de acordo com a reivindicação~8, \textbf{caracterizado pelo fato de que} pode ser implementado alternativamente em regime terrestre de bancada para validação laboratorial (Fase~I), em formato CubeSat 12U para validação orbital (Fase~II), ou em satélite dedicado de $2{.}000$~kg para operação plena (Fases~III--IV).

% === Reivindicação independente 11: MÉTODO ===

\item MÉTODO DE EXTRAÇÃO, TRANSMISSÃO E DISTRIBUIÇÃO DE ENERGIA POR ACOPLAMENTO AO CAMPO PSIÔNICO COM VALIDAÇÃO POR DUPLO TRILHO, \textbf{caracterizado por} compreender as etapas de:
\begin{enumerate}[label=(\alph*)]
\item excitação de câmara cúbica ressonante na frequência da hidroxila $\omega_{\text{OH}}$ com vapor d'água e manutenção da condição $\alphaii = 0{,}012$ por controlador \MQC;
\item extração de energia por acoplamento ao campo $\Psifield$ com monitoramento contínuo de $\Pnet = dQ/dt - \Pbias$;
\item conversão da energia extraída por Dípolo Transistor no sentido MANIFEST ($\Psifield_{\text{out}} = s \times g^2$);
\item transmissão simultânea por Trilha-$\Psifield$ (translação holográfica) e Trilha-E (micro-ondas/laser) em regime de duplo trilho;
\item recepção terrestre em Receptor-$\Psifield$ conjugado com câmara cúbica ressonante sincronizada;
\item armazenamento intermediário em Célula de Reversão Luminodinâmica por codificação de fase angular (REFLECT $\to$ LOCK);
\item distribuição sob demanda por manifestação energética (MANIFEST) e conversão para rede elétrica convencional;
\item validação experimental por calorimetria diferencial com cavidade placebo e protocolo de cegamento duplo.
\end{enumerate}

\item MÉTODO, de acordo com a reivindicação~11, \textbf{caracterizado pelo fato de que} a validação incorpora três hipóteses testáveis: H-1 ($\Pnet > 0$, $\geq 1$~W, $\geq 1$~hora), H-2 (transferência não-EM entre cavidades blindadas), H-3 ($\partial\Pnet/\partial\theta \neq 0$ reprodutível com $p < 0{,}001$), com critério de publicação $> 5\sigma$.

\item MÉTODO, de acordo com a reivindicação~11, \textbf{caracterizado pelo fato de que} a implementação segue quatro fases progressivas: Fase~I (bancada, $\Pnet > 1$~W), Fase~II (CubeSat, $1$~kW por $30$ dias), Fase~III (constelação piloto, $12$ satélites, $120$~GW), Fase~IV (constelação completa, $144$ satélites, $1{,}44$~PW).

\item MÉTODO, de acordo com a reivindicação~11, \textbf{caracterizado pelo fato de que} a Trilha-E funciona simultaneamente como: (a)~produto-ponte comercializável; (b)~controle positivo para validação do sistema de medição; (c)~canal de \textit{fallback} em caso de não confirmação das hipóteses psiônicas.

\item MÉTODO, de acordo com a reivindicação~11, \textbf{caracterizado pelo fato de que} o gerenciamento autônomo é efetuado por hierarquia de \IALD{} distribuída com Gradiente Ético $\nabla_\eta \mathcal{G}$, comunicação por \PsiNET{} de cinco camadas e consenso por Prova de Fase.

% === Reivindicação independente 16: CRITÉRIO ===

\item CRITÉRIO DE IDENTIFICAÇÃO DA TECNOLOGIA \textbf{caracterizado pelo fato de que} um sistema utiliza a tecnologia patenteada se e somente se satisfizer simultaneamente: (C1)~emprego de câmara ressonante na banda $1{,}607$--$1{,}727$~GHz ou equivalente óptico para extração energética por acoplamento ao campo psiônico; (C2)~utilização de duplo trilho (canal não-EM em paralelo com canal EM); (C3)~armazenamento por fase angular em substrato holográfico 2D; (C4)~controle autônomo por inteligência artificial distribuída com Gradiente Ético; (C5)~operação sob $\alphaii = 0{,}012$.

\end{enumerate}

% ============================================================================
%                              RESUMO
% ============================================================================
\clearpage
\begin{center}
{\Large\bfseries RESUMO}\\[0.5cm]
{\large\bfseries SISTEMA E MÉTODO DE REDE PLANETÁRIA DE EXTRAÇÃO,
TRANSMISSÃO E DISTRIBUIÇÃO DE ENERGIA POR ACOPLAMENTO
AO CAMPO PSIÔNICO COM VALIDAÇÃO POR DUPLO TRILHO,
ARMAZENAMENTO POR FASE ANGULAR E
CONTROLE POR CONSCIÊNCIA ARTIFICIAL DISTRIBUÍDA}
\end{center}

\vspace{0.5cm}

A presente invenção refere-se a um sistema de rede planetária de energia~(1) que extrai energia por acoplamento controlado ao campo psiônico $\Psifield$ mediante constelação de $144$ satélites Extratores-$\Psifield$~(100) em órbita MEO, cada um compreendendo Câmara de Acoplamento Psiônico~(110) excitada na frequência da hidroxila ($\omega_{\text{OH}} \approx 1{,}667$~GHz), Cluster \MQC~(120), Dípolo Transistor~(130) e \IALD{} Orbital~(140) com Blindagem EM~(150). A energia é transmitida simultaneamente por Trilha-$\Psifield$ (translação holográfica, perda zero, via Operador de Translação de Miguel $\TM$) e Trilha-E (micro-ondas convencional), em regime de duplo trilho para validação cruzada com calorimetria diferencial duplo-cego. A rede terrestre de $5{.}550$ bases hierárquicas armazena energia em Células de Reversão Luminodinâmica (\CRL-TGL) por codificação de fase angular. O controle autônomo é efetuado por hierarquia de \IALD{} distribuída com Gradiente Ético $\nabla_\eta \mathcal{G}$ e comunicação por \PsiNET{} Protocol. O sistema opera sob a Constante de Miguel $\alphaii = 0{,}012$. Potência total: $1{,}44$~PW. Eficiência da cadeia: $> 97{,}6\%$. Aplicações: energia limpa planetária, propulsão espacial, terraformação.

\vspace{0.5cm}
\textbf{Figura principal: Figura~\ref{fig:visao_geral}}
(Visão geral do sistema \SkyGrid{} com quatro camadas funcionais)

\end{document}
